% =============================================================================
%  Capítulo 6 -- Discusión
% =============================================================================
\chapter{Discusión}
\label{cap:discusion}

\section{Implicaciones sociotécnicas de los resultados}
\label{sec:implicaciones}

[Articular los hallazgos cuantitativos del Capítulo~\ref{cap:resultados} y
los valores del \INS del Capítulo~\ref{cap:ins} con la discusión teórica del
Capítulo~\ref{cap:marco}. ¿Qué dicen estos resultados sobre la traslación de
la discriminación lingüística humana al ámbito algorítmico? ¿Qué dicen sobre
las prácticas actuales de evaluación de sistemas de \PLN desplegados en
contextos hispanohablantes?]

\subsection{Implicaciones para el desarrollo de sistemas}

[Lecciones aplicables a equipos que diseñan y despliegan clasificadores de
texto en contextos como el filtrado de currículums.]

\subsection{Implicaciones para la evaluación de equidad}

[¿Cuándo conviene incorporar el \INS o variantes suyas en las baterías de
evaluación de un sistema de \PLN?]

\subsection{Implicaciones para la discusión pública sobre algoritmos}

[Relación de los resultados con debates más amplios sobre transparencia,
rendición de cuentas y regulación de sistemas algorítmicos.]

\section{Limitaciones metodológicas del estudio}
\label{sec:limitaciones}

\subsection{Limitaciones del componente social}

[Sobre la muestra de la encuesta, su representatividad, los sesgos del
instrumento.]

\subsection{Limitaciones del \textit{benchmark} controlado}

[Naturaleza simulada de los textos; cobertura limitada de los tipos de
variación; subjetividad en el control de equivalencia informativa.]

\subsection{Limitaciones del indicador propuesto}

[El \INS como indicador preliminar; supuestos implícitos; condiciones bajo
las cuales podría sobreestimar o subestimar la neutralidad efectiva del
sistema.]

\subsection{Limitaciones de la generalización}

[Acotaciones explícitas sobre la transferibilidad de los hallazgos a otras
lenguas, contextos y dominios de aplicación.]
